% ============================================================
%  2. fejezet, Irodalmi háttér
% ============================================================
\chapter{Irodalmi háttér}
\label{ch:irodalom}

\section{Kalman-szűrő és általánosításai}
\label{sec:kalman}

A Kalman-szűrő~\cite{kalman1960new} minimális varianciájú, torzítatlan,
lineáris becslő olyan rendszerekre, amelyek $x_t \in \mathbb{R}^n$ állapota
Gauss-zajjal perturbált lineáris dinamikai modell szerint fejlődik:
\begin{align}
    x_t &= F x_{t-1} + w_t, & w_t &\sim \mathcal{N}(0, Q), \label{eq:kf_process}\\
    y_t &= H x_t   + v_t,   & v_t &\sim \mathcal{N}(0, R), \label{eq:kf_measure}
\end{align}
ahol $F$ az állapotátmeneti mátrix, $H$ a mérési mátrix,
$Q \in \mathbb{R}^{n \times n}$ a folyamatzaj-kovariancia, és
$R \in \mathbb{R}^{m \times m}$ a mérési-zaj-kovariancia.
A szűrő két felváltva ismétlődő fázisban működik.

Az \emph{előrejelzési lépésben} (időfrissítés), mielőtt a $y_t$ mérés
megérkezne, a szűrő az állapotbecslést és a hibakovariancát előrevetíti:
\begin{align}
    \hat{x}_{t|t-1} &= F \hat{x}_{t-1|t-1}, \label{eq:kf_predict_x}\\
    P_{t|t-1}       &= F P_{t-1|t-1} F^\top + Q, \label{eq:kf_predict_P}
\end{align}
ahol $P_{t|t-1} \in \mathbb{R}^{n \times n}$ az előrejelzett hibakováriancia.
Nagy $P_{t|t-1}$ értéke magas előrejelzési bizonytalanságot jelez. Nagy $Q$ azt jelenti,
hogy az állapot gyorsan és kiszámíthatatlanul változik.

A \emph{korrekciós lépésben} (mérésfrissítés), amint a $y_t$ megfigyelés
megérkezik, a szűrő az \emph{innováció} $e_t = y_t - H\hat{x}_{t|t-1}$
alapján korrigálja becslését, a Kalman-gain $K_t$ által súlyozva:
\begin{align}
    K_t           &= P_{t|t-1} H^\top \bigl(H P_{t|t-1} H^\top + R\bigr)^{-1},
                     \label{eq:kalman_gain}\\
    \hat{x}_{t|t} &= \hat{x}_{t|t-1} + K_t e_t, \label{eq:kalman_update}\\
    P_{t|t}       &= (I - K_t H)\,P_{t|t-1}. \label{eq:kf_update_P}
\end{align}
A~\eqref{eq:kalman_gain}. egyenlet a zaj--jel-egyensúly alapvető összefüggését tükrözi.
Ha az $R$ mérési zaj nagy az előrejelzett $H P_{t|t-1} H^\top$
kovarianciához képest, a nevező dominál és $K_t \to 0$: a szűrő nem bízik
a mérésben, és az előzmény-állapotra támaszkodik. Ha $R \to 0$, akkor
$K_t \to H^{-1}$, és a megfigyelés majdnem teljesen felülírja az állapotot.

A skaláris esetben ($F = H = 1$) a gain egyszerűsödik:
\begin{equation}
    K_t = \frac{P_{t|t-1}}{P_{t|t-1} + R}.
    \label{eq:kalman_gain_scalar}
\end{equation}
Stacionárius feltételek mellett a $P_{t|t-1}$ hibakováriancia egy fix
$\bar{P}$ értékre konvergál, és az egyensúlyi gain kizárólag a $Q/R$
aránytól függ: magas folyamatzaj ($Q$) a $K$-t $1$ felé tolja (bízzunk
a mérésben), alacsony $Q/R$ arány $K$-t $0$ felé tolja (bízzunk
a modellben). (\ref{sec:kalman_gain}. szakasz).

\subsection{Tanult Kalman-szűrők}

A neurális Kalman-irodalomban a $F, H, Q, R$ mátrixokat neurális hálózatokkal
paraméterezik~\cite{krishnan2015deep,fraccaro2017disentangled}.
A \emph{Deep Kalman Filter}~\cite{krishnan2015deep} nemlineáris átmenetet
enged meg, a \emph{Kalman Variational Autoencoder}~\cite{fraccaro2017disentangled}
a latens állapotok szétválasztott reprezentációját tanítja.
A KCMamba ezeken az elveken épít, de nem autoencoderként, hanem
szekvencia-előrejelzőként.

\section{Állapottér-modellek}
\label{sec:ssm}

\subsection{S4}

Az S4 (Structured State Spaces for Sequences)~\cite{gu2021s4} egy
folytonos-idejű állapottér-modellből indul ki, amelyben egy $d$-dimenziós
bemeneti sorozatot egy $N$-dimenziós látens dinamikán keresztül modellez:
\begin{align}
    h'(t) &= A h(t) + B u(t), \\
    y(t) &= C h(t),
\end{align}
ahol $A \in \mathbb{R}^{N \times N}$ az állapotátmeneti mátrix, $B \in \mathbb{R}^{N \times d}$
a bemeneti projekció, $C \in \mathbb{R}^{d \times N}$ a kimeneti projekció,
és $h(t) \in \mathbb{R}^N$ a látens állapot.
Diszkretizálva egy $\Delta$ időlépéssel:
\begin{equation}
    h_t = \bar{A}\, h_{t-1} + \bar{B}\, u_t, \quad y_t = C\, h_t,
\end{equation}
ahol $\bar{A} = e^{\Delta A}$ és $\bar{B} \approx (\Delta A)^{-1}(e^{\Delta A} - I) \Delta B$.
Az S4 kulcsinnovációja az $A$ mátrix HiPPO-inicializációja, amely
az állapotnak polinomiális bázisú hosszú távú memóriát biztosít.
A konvolúciós nézet, $\bar{K} = (C\bar{B}, C\bar{A}\bar{B}, \ldots, C\bar{A}^{T-1}\bar{B})$
kernel, lehetővé teszi, hogy a tanítás $O(T \log T)$-ben
gyors Fourier-transzformációval (FFT) végezhető, míg az inferencia a rekurrens formában $O(1)$ memóriával
fut lépésenként.

\subsection{Mamba}

A Mamba~\cite{gu2023mamba} az S4 alapjaira építve bevezeti az
\emph{input-függő szelektivitást}: míg az S4-ben az $A$, $B$, $C$
mátrixok rögzítettek a teljes szekvenciára, a Mamba ezeket az aktuális
bemenettől függővé teszi. Ez a szelektivitás a lényege: a modell
minden időlépésben eldönti, mit őrizzen meg az állapotban és mit
felejtsen el. A bemeneti projekciós mátrix
$B(u_t) \in \mathbb{R}^{N \times d}$, a kimeneti projekció
$C(u_t) \in \mathbb{R}^{d \times N}$ és az időléptékező
$\Delta(u_t) \in \mathbb{R}^d$ mindegyike az aktuális
$u_t \in \mathbb{R}^d$ bemeneti vektortól függ:
\begin{equation}
    h_t = \overline{A}(u_t)\, h_{t-1} + \overline{B}(u_t)\, u_t,
    \quad y_t = C(u_t)\, h_t,
    \label{eq:mamba_recurrence}
\end{equation}
ahol $h_t \in \mathbb{R}^N$ a látens állapotvektor ($N$ az állapotméret),
$\overline{A}(u_t) = e^{\Delta(u_t) A}$ és
$\overline{B}(u_t) \approx \Delta(u_t) B$ a $\Delta$ alapján diszkretizált
mátrixok.
A modell így szelektíven dönti el, hogy a bemeneti sorozat melyik részét
tárolja tartósan az állapotban, és melyiket hagyja figyelmen kívül.
Hardver-barát párhuzamos prefix-scan implementációval $O(T)$ memória- és
időkomplexitást ér el.

\subsection{Fair Mamba}

Az összehasonlítási alap, amelyet \emph{Fair Mamba} névvel illetünk,
egy csatornánkénti (\emph{channel-wise}) Mamba, amelyet a KCMamba-éval
azonos hiperparaméterekkel és tanítási protokoll szerint tanítunk.
A \emph{fair} jelző arra utal, hogy a paraméterbudget és az adatmennyiség
azonos, így az eredménybeli különbség kizárólag az architektúrából fakad.

\section{Hosszú távú idősor-előrejelzés}
\label{sec:ltsf}

\subsection{Problémaformulázás}

A hosszú távú idősor-előrejelzés (LTSF) egy
$\mathbf{X} = (x_1, \ldots, x_L) \in \mathbb{R}^{L \times d}$
bemeneti ablakot tekint, amelynek $L$ historikus lépése $d$ csatornán van megadva,
és $\hat{\mathbf{X}} = (\hat{x}_{L+1}, \ldots, \hat{x}_{L+H}) \in \mathbb{R}^{H \times d}$
előrejelzését kívánja adni a következő $H$ lépésre.
A cél a normált átlagos négyzetes hiba minimalizálása:
\begin{equation}
    \mathrm{MSE} = \frac{1}{H d}
        \sum_{h=1}^{H} \sum_{j=1}^{d}
        \frac{(x_{L+h,j} - \hat{x}_{L+h,j})^2}{\mathrm{Var}(x_{\cdot,j})},
    \label{eq:mse}
\end{equation}
ahol a csatornánkénti varianciával való normálás különböző skálájú
adatkészletek eredményeit összehasonlíthatóvá teszi.
A \emph{hosszú távú} jelző a $H \geq L$, illetve $H$ és $L$ hasonló nagyságrendű esetét jelöli:
az ilyen feladatban a modellnek a visszatekintési ablakán jócskán túlra kell előrejelezni,
ez pedig lényegesen nehezebb, mint a rövid távú eset, ahol $H \ll L$.

\subsection{Referenciamodellek}

\emph{Informer}~\cite{zhou2021informer} volt az első, kifejezetten hosszú
távú előrejelzésre tervezett Transformer-változat. A standard önfigyelés
$O(L^2)$ idő- és memóriakomplexitása hosszú bemeneti ablakoknál
elfogadhatatlanul magas. Az Informer ezt \emph{ProbSparse attention}\textrm{-}tel
orvosolja: KL-divergencia-becslés alapján azonosítja az $O(\log L)$
leginformatívabb lekérdező pozíciót, ezzel $O(L \log L)$-re csökkenti
a komplexitást. Disztillációs művelet felezi a szekvenciahosszat az
enkoderblokkok között, így több ezer lépéses visszatekintési ablakok is
kezelhetők.

\emph{Autoformer}~\cite{wu2021autoformer} minden sorozatot egy
trendkomponensre (mozgóátlag-kernellel becsülve) és egy szezonális
maradéktagra bont, és a Transformert kizárólag a maradéktagokra alkalmazza.
Pontszorzat-figyelem helyett \emph{Auto-Korrelációs} mechanizmust
használ: megkeresi azt a $\tau$ késleltetést, amely maximalizálja
a lekérdező jel autokorrelációját, majd ezeken a késleltetési pozíciókon
összesíti az értékeket. A mechanizmus FFT-vel $O(L \log L)$ komplexitású,
és periodikus jelekre illeszkedő induktív torzítással rendelkezik.

\emph{DLinear}~\cite{zeng2023transformers} elveti az architekturális
bonyolultságot, és az előrejelzést minden visszatekintési ablakra
egyetlen lineáris réteggel állítja elő, opcionálisan trend-szezonális
felbontás után. Egyszerűsége ellenére a DLinear számos Transformer-változattal
azonos szinten teljesített, sőt azokat felülmúlta szabványos benchmarkokon,
bizonyítva, hogy az architekturális bonyolultság nem feltétele a jó
előrejelzési teljesítménynek. Referenciamodellként releváns marad, mert
minden új modellnek igazolnia kell, hogy a magasabb számítási ráfordítás
indokolt ehhez az egyszerű alapmodellhez képest.

\subsection{Adatkészletek}
\label{sec:datasets}

Az Exchange-Rate adatkészlet~\cite{lai2018modeling} nyolc valutaárfolyamot
tartalmaz, amelyeket 1990 és 2016 között napi szinten vettek fel,
csatornánként 7\,588 időlépéssel. Az árfolyam-adatok dinamikája majdnem egységgyök-jellegű: minden napi
változás véletlenszerű bolyongáshoz közelít, szinte elhanyagolható
vissza-átlagolással. Ez kemény feladatot jelent az autoregresszív modelleknek,
és természetes tesztkörnyezetet jelöl a Kalman-struktúra számára,
ahol a jel-zaj arány eleve alacsony.

Az ETTh1 és ETTh2 adatkészletek~\cite{zhou2021informer} kínai
transzformátorállomásokból gyűjtöttek óránkénti villamosenergia-terhelést
és olajhőmérsékletet két éven keresztül (17\,420 lépés, $d=7$ csatorna).
Az ETTh1 egyértelmű napi periodicitást mutat, az ETTh2 zajosabb és
szabálytalanabb. A két adatkészlet eltérő nehézségű tesztkörnyezetet
nyújt a zajrobusztusság kiértékeléséhez.

\section{Zajrobusztusság}
\label{sec:robustness}

Az idősor-modellek zajérzékenységének szisztematikus elemzése az
LTSF-irodalomban aláreprezentált terület.
A PatchTST~\cite{nie2022patchtst} patch-tokenizálással csökkenti a
szomszédos mintavételezési torzítást. Ennek közvetett zajcsökkentő hatása
van, bár nem ez az eljárás fő célkitűzése.

A jelen munkában az értékelési protokoll: $\tilde{x}_t = x_t + \varepsilon_t$,
ahol $x_t \in \mathbb{R}^d$ az eredeti idősorjel, $\varepsilon_t \sim \mathcal{N}(0, \sigma^2 I)$
additív Gauss-zaj ($\sigma > 0$ a szórás, $I \in \mathbb{R}^{d \times d}$
az egységmátrix), és $\sigma \in \{0, 1, 3, 5\}$ a vizsgált zajszintek halmaza.
A tanító- és teszthalmazt egyaránt azonos $\sigma$-val zajosítjuk,
tehát nem domain-shift-ről, hanem egyenletes adatminőség-romlásról van szó.
